\documentclass[10pt,conference]{IEEEtran}
\usepackage{cite}
\usepackage{amsmath}
\usepackage[T1]{fontenc}
\usepackage{graphicx}
\usepackage{url}
\usepackage{hyperref}
\usepackage{booktabs}
\usepackage{tabularx}
\usepackage{array}
% Author details preserved from paper.tex; excluded from anonymous rendering:
% Vitor Raposo, Ontario Tech University, Oshawa, Ontario, Canada
% vitor.raposo@ontariotechu.net
% Sanaa Alwidian, Ontario Tech University, Oshawa, Ontario, Canada
% sanaa.alwidian@ontariotechu.ca
% Study repository:
% https://github.com/MyLittleFoxxie/MESA-meta-evaluation-standard-for-ai

\begin{document}
\title{Making Artificial Intelligence Benchmark Review More Approachable}
\author{Anonymous authors}
\maketitle
\begin{abstract}
\textbf{Background:} Reviewing an AI benchmark requires more than reading its scores. A reader needs to understand how the test works, what was checked, and which conclusions the available information supports. Existing review frameworks provide criteria, but applying those criteria can require knowledge of measurement.
\textbf{Objective:} We present the Meta-Evaluation Standard for AI Benchmarks (MESA), a questionnaire for reviewing how benchmarks are documented, checked and used to support claims. Its 25 questions adapt the Test Review Model of the European Federation of Psychologists' Associations. Explanations and five answer choices aim to help readers without measurement training apply the questions. We examine what this adaptation preserves and how its supporting material affects answers.
\textbf{Method:} We describe the adaptation in relation to EFPA, OECD and AI evaluation research. Two model studies are planned: questions with and without explanations, and original source passages compared with survey summaries. Reference answers distinguish information available in each presentation and allow documented alternatives where a single answer cannot be established.
\textbf{Results:} [Pending: report completed model sessions, reference-answer comparisons, interpretation problems and the effects of summaries. The proposed experiments have not been conducted.]
\textbf{Conclusion:} The adaptation and study design make MESA's scope and assumptions inspectable. Whether explanations support answers and whether summaries preserve relevant information remain empirical questions. Human approachability will be examined separately through the existing participant survey.
\end{abstract}
\begin{IEEEkeywords}
AI benchmarks, test review, questionnaire, approachability, EFPA-TRM, score interpretation
\end{IEEEkeywords}

\noindent\textbf{Draft status, 29 October 2026.} This revision contains the proposed study method and explicit result placeholders. Study materials and reference answers require completion and checking before execution. No findings from the earlier model study are presented as results of these experiments.

\section{Introduction}
An engineer choosing an AI component may begin by comparing benchmark scores. Before relying on a difference between systems, however, the engineer needs to understand how those scores were obtained. Did the test resemble the proposed use? Did the systems receive comparable resources? Would the difference remain if the test were repeated? Answering these questions requires examining both the test and the conclusions drawn from it. Guidance from psychological and educational testing provides a basis for connecting score interpretations and uses to the information that supports them \cite{aera_apa_ncme_standards_2014}.

People who need to ask these questions may have no training in measurement. This creates a need for approachability in benchmark review: helping readers understand review questions, find relevant information, explain their answers and recognise what remains unknown. A shorter questionnaire may provide a starting point, but its length alone cannot show whether readers can use it. Some assessments will still require expertise, and faster completion is not the sole purpose.

Several sources already help reviewers examine tests. The European Federation of Psychologists' Associations Test Review Model (EFPA-TRM) offers 90 pages of forms and guidance for expert reviewers \cite{efpa_test_review_model_2025}. OECD describes facets of AI evaluation instruments and reports that some classifications required further reading and discussion \cite{cohn_hernandez_orallo_oecd_2023}. BetterBench and Evidence-Centered Benchmark Design (ECBD) provide criteria and guidance for AI benchmarks \cite{reuel_betterbench,liu_ecbd_2024}. Building on this work raises a further question: how can readers apply selected criteria through a shorter form while retaining distinctions that matter to the answer?

To investigate this question, we present the Meta-Evaluation Standard for AI Benchmarks (MESA), a questionnaire adapted from EFPA-TRM. MESA organises a review around 25 questions about how a benchmark is described, developed, scored, checked and used to support claims. A reviewer reads the available benchmark information and answers each question using one of five choices: Yes, Partly, No, Not sure or Not applicable. The questions are arranged in six domains, with explanations to help readers understand what to look for. The study materials add definitions, benchmark descriptions and a completed example. Together, these materials put MESA's aim of approachability into practice; whether they help readers remains to be examined. Answers remain separate rather than being combined into a benchmark quality score. The word Standard in the name does not imply approval by a standards organisation.

Moving from review guidance to a questionnaire introduces two choices about how information is presented. The first is how to explain each question. An explanation might help a reader distinguish a missing report from a reported problem. The second is how to present the benchmark documentation. A summary might make a fact easier to find, while also omitting a qualification or suggesting an answer. We examine these choices separately to understand how each affects the review.

We begin by tracing what MESA keeps, changes and leaves out of its source guidance. We then describe how reference answers can be linked to the information supplied to reviewers, and propose two model studies of the explanations and summaries. Together, these steps address three research questions:
\begin{enumerate}
\item \textbf{RQ1:} Which review concerns does MESA keep, change or leave out?
\item \textbf{RQ2:} Do MESA's explanations help AI reviewers answer its questions using the information supplied?
\item \textbf{RQ3:} How do AI reviewers' answers change when original passages are replaced with the survey summaries?
\end{enumerate}
RQ1 is addressed through the adaptation analysis. Studies 1 and 2 will address RQ2 and RQ3. They test actual model responses, not simulated people. The existing human survey is Study 3 in the wider research programme and is reserved for a later extension. Its completion is not a requirement for the CHASE model study.

\section{Background and Related Work}
\subsection{Reviewing a test and interpreting its scores}
The AERA/APA/NCME Standards locate validity in the evidence supporting interpretations and uses of scores \cite{aera_apa_ncme_standards_2014}. For example, correctly counting solved problems does not establish that a total measures every part of reasoning. Reliability concerns consistency and precision under specified conditions; fairness depends on the people or systems, conditions and uses under consideration. These concepts cannot be replaced by a single statement that a benchmark is good or bad.

EFPA-TRM separates describing a test from evaluating its rationale, materials, comparisons, reliability, validity, fair use and reports \cite{efpa_test_review_model_2025}. Its explanations and rating anchors support expert review. Transfer to AI changes the object of some checks: marking may be performed by software or another model; exposure may occur during training; and test conditions include prompts, tools and model settings. MESA's adaptation retains selected functions while documenting these changes.

\subsection{Frameworks for reviewing benchmarks}
OECD Chapter 6 groups 18 facets under Validity, Consistency and Fairness \cite{cohn_hernandez_orallo_oecd_2023}. Some classify an instrument rather than rank its quality. Judgeability identifies manual, automated or mixed scoring; it does not say whether scoring is correct. The study reports that Coverage and Specificity sometimes required related literature beyond the instrument's specification. This motivates support for interpretation, without implying that every facet is difficult.

Other frameworks move from describing a benchmark to examining its development and use. BetterBench covers documentation, implementation, reproduction and maintenance \cite{reuel_betterbench}. ECBD includes a worksheet asking designers to describe, justify and support their decisions \cite{liu_ecbd_2024}. These approaches already provide structured guidance. MESA builds on that context by examining how selected review functions are translated into questions and how the accompanying material affects their application.

\subsection{Problems that review questions need to expose}
Measuring what Matters examines links among the ability of interest, tasks, metrics and claims \cite{bean_measuring_what_matters}. The Leaderboard Illusion shows how participation and reporting practices affect the interpretation of rankings in its studied platform \cite{singh_leaderboard_illusion}. These sources motivate questions about task coverage, comparison conditions and public conclusions. They do not establish that one questionnaire resolves those problems.

A Definition of AGI provides a proposed account of broad cognitive abilities \cite{hendrycks_definition_agi}. It is relevant when a benchmark result is used to support a general intelligence claim. It is not a universal basis for evaluating every task benchmark, nor a reason to convert MESA answers into an intelligence or quality score.

\subsection{Scope of the literature review}
A supporting review covers publications from 28 October 2023 through 28 October 2026, with the 2014 Standards as a foundation exception. It records searches of OpenAlex, Crossref and arXiv, supplemented by institutional sources and citation searching. Search caps, selected backward citations, access failures and unresolved texts limit coverage. Its source types include standards, proposals and empirical studies; their recommendations and findings are not interchangeable. Assistant-prepared inclusion decisions and gap claims await author review. This draft treats approachability as a research question, not as a demonstrated absence of prior solutions.

\section{The Questionnaire and Its Basis}
\subsection{Questions, explanations and choices}
MESA's current questions are preserved in this revision. The preparation snapshot is participant packet 2.9, identified as \texttt{mesa-guided-arc-designqa-v2-9}. The first five questions concern whether descriptions are documented; later questions concern reasons, checks and the interpretation of results. Table~\ref{tab:domains} summarises the six domains. Full wording, explanations, definitions and source references are retained in the supporting materials.

\begin{table}[t]
\caption{The existing MESA domains.}
\label{tab:domains}
\small
\begin{tabularx}{\columnwidth}{@{}p{0.08\columnwidth}X@{}}
\toprule
Q & Domain and concern \\
\midrule
1--5 & Description: identity, tasks, conditions and scoring. \\
6--9 & Purpose and development: reasons, checking problems and coverage. \\
10--13 & Scoring and comparisons: points, combined scores, marking and reference results. \\
14--16 & Consistency of results: repeated runs, changed conditions and variation. \\
17--21 & Meaning and fairness: shortcuts, exposure, training, claims and comparisons. \\
22--25 & Use and reporting: inspection, reuse, updates and public claims. \\
\bottomrule
\end{tabularx}
\end{table}

Yes means the information supports the whole condition; Partly means it supports some parts. No requires information showing that the condition is not met. Not sure records missing, unclear or insufficient information. Not applicable means that the issue does not arise for the benchmark or intended use. These choices describe answers to individual questions, not numerical levels of benchmark quality.

The wording of the question determines what is being assessed. For a question about whether a check occurred, finding errors does not change the fact that checking occurred. For a question about whether documentation is available, an explicitly incomplete description concerns reporting, not necessarily the underlying procedure. Likewise, a reviewer can recognise a stated rationale without establishing that the rationale is empirically supported.

\subsection{What the questionnaire keeps and leaves out}
The EFPA correspondence is an explanation of the design, not a recovered history of every development decision or evidence of endorsement. For each review function, the analysis records its source, related MESA question, change in meaning and remaining limit. Table~\ref{tab:adaptation} gives examples relevant to RQ1. Matching a heading does not establish equivalent coverage, so no overall preservation percentage is calculated.

\begin{table*}[t]
\caption{Examples of adaptation decisions and remaining limits.}
\label{tab:adaptation}
\small
\begin{tabularx}{\textwidth}{@{}p{0.15\textwidth}p{0.05\textwidth}XX@{}}
\toprule
Review function & Q & MESA treatment & What the answer does not establish \\
\midrule
Describe the test & 1--5 & Combine identity, tasks, conditions and scoring into documentation questions. & A complete technical specification or adequate testing practice. \\
Check content and marking & 8, 12 & Keep problems and supplied answers separate from the method marking new responses. & That checking found no errors or covered every failure mode. \\
Examine coverage & 7, 9 & Ask about selection and coverage relative to the claimed ability. & Detailed analysis of how individual items relate to the ability. \\
Interpret comparisons & 13 & Replace an assumed human reference population with justified comparison points. & Representativeness of every model cohort or a validated decision threshold. \\
Examine consistency & 14--16 & Separate repeated conditions, changed conditions and estimates of variation. & All forms of reliability or equivalence between test versions. \\
Address AI conditions & 18, 19, 24 & Ask about training exposure, submission rules and changes affecting scores. & Knowledge of proprietary training data or a full audit of historical results. \\
\bottomrule
\end{tabularx}
\end{table*}

For example, Q8 asks whether problems and their supplied answers were checked during development, while Q12 asks whether marking was checked. An audit of reference answers does not automatically test a model used as a judge. Q14 concerns repeated testing under the same setup; Q15 concerns changes to that setup. A range across changed prompts cannot establish variation under identical conditions. Q13 retains the need for justified comparisons but omits separate EFPA criteria for population norms and decision cutoffs.

These distinctions also expose the cost of combining concerns. MESA does not separately reproduce all EFPA criteria for internal consistency, item models, human test-user qualifications or reference populations. Further analysis is needed when the intended claim depends on those functions.

\subsection{Connections to other review frameworks}
OECD terminology helps explain correspondences after MESA's development. Coverage relates to Q7 and Q9; Reliability to Q14 and Q16; Specificity to alternative ways of earning points in Q10 and Q17. These are conceptual relations, not interchangeable definitions. OECD's Specificity includes confounding by other abilities; its use of contamination should not be reduced to prior training exposure.

Other correspondences are incomplete. Q5 describes scoring, whereas OECD's Judgeability classifies who or what performs it. Q15 examines changed conditions without fully classifying Adjustability or Variation. OECD's Progression concerns learning during the test and has no explicit MESA item; it is not the release maintenance addressed by Q24. The supporting map records all 18 facets, including weak or absent correspondences. The six MESA domains remain unchanged.

\section{Reference Answers and Study Materials}
\subsection{Benchmark scope and source versions}
To examine how the questions work with benchmark documentation, we use two different settings: the grid problems of ARC-AGI-2 and the engineering documentation of DesignQA \cite{chollet_arc_agi_2_2025,doris_designqa_2024}. We selected these benchmarks for the study; they do not represent all benchmark types. The reading packet identifies their versions, intended uses and source inspection dates. Those source versions must be fixed before execution so that later changes do not alter the comparison.

Humanity's Last Exam (HLE) remains a completed teaching example \cite{phan_etal_hle_2025}. It is supplied in every condition and is not scored as an unseen task. The model is asked to review descriptions, not solve benchmark problems or run the benchmark software.

\subsection{Two reference answers for each real question}
Comparing reviewers' answers requires first deciding what the supplied information supports. For each of the 25 questions and both benchmarks, a reference record will link a proposed answer to a source version and a passage or code location. It will also explain the answer and record acceptable alternatives. Because a summary may contain different information from the original passages, each question needs a reference for each reading.

Suppose an original passage reports repeated testing but the summary omits it. Yes may be supported by the original passage and Not sure by the summary. The latter answer should not be marked wrong for failing to recover an omitted fact. Table~\ref{tab:example} illustrates this distinction with fictional statements; it is not a finding about either benchmark.

\begin{table}[t]
\caption{Illustration for Q14: was the same setup tested repeatedly?}
\label{tab:example}
\small
\begin{tabularx}{\columnwidth}{@{}Xp{0.18\columnwidth}@{}}
\toprule
Information supplied & Reference \\
\midrule
Five runs used the same model and settings to measure score variation. & Yes \\
The authors state that only one run was performed. & No \\
The description does not say whether runs were repeated. & Not sure \\
\bottomrule
\end{tabularx}
\end{table}

Some questions permit a factual check, while others require an assessment of adequacy or claim scope. Reference records therefore distinguish assessable, contested and unresolved items. A contested item may have several defensible choices with different reasons. Main label-match counts will use the subset with a settled reference and report its denominator; all other items remain visible. Sensitivity analyses will show the effect of documented alternatives. These decisions must precede new responses, not be selected to favour a condition.

\textbf{Preparation status.} The reference register has 50 entries, but author checking and source-passage reference answers are pending. The packet's statements are not substitutes for checking the original sources. No reference register is yet frozen for scoring.

\subsection{Fictional cases and checking summaries}
The design specifies 24 question--description cases in six groups: missing information versus explicit absence; checking versus outcome; repeated versus changed conditions; scoring description versus marking checks; available materials versus reproduction; and tested tasks versus broader claims. Each group has four cases targeting a current MESA question. This remains a focused check rather than a comprehensive test. Choices are used only where the question supports them, without forcing equal category frequencies. These cases assess selected distinctions, not all 25 questions equally.

The original-passage packet will use a documented selection procedure across the same source set as the summary. Its audit records retained facts, omitted qualifications, added interpretations and direct answer cues. Original excerpts are also selected by researchers and are not inherently unbiased. Software excerpts must be distinguished from a researcher's interpretation of the software. Changes in information amount and readability will be recorded as possible explanations for answer differences.

The author must check case labels and reference records before execution. Additional checking by someone other than the preparer is desirable, but will be reported only if performed. Existing model answers and optional human comments are not used to construct a convenient answer key. The preparation record must disclose whether any human responses were already available to analysts.

\section{Study 1: Questions and Explanations}
\subsection{Conditions and tasks}
The first study asks whether explanations help models apply the questions. Condition A supplies the questions and common answer instructions. Condition B adds the question explanations and associated definitions. Both receive the same choice meanings, response format and completed HLE example. Keeping these materials in both versions lets us examine what the additional explanations contribute when a teaching example is already available.

Fictional cases and real benchmark reviews are separate tasks. For the real task, A and B both receive the existing survey summaries and answer all 25 questions for each benchmark. Each answer contains a choice, a short reason and a supporting passage identifier or quotation where applicable. When information is absent, the response can identify what is missing instead of inventing a quotation.

\subsection{How the studies will run}
The selected reviewer models are \texttt{gpt-6-luna} and \texttt{gpt-6.1-terra}, with low reasoning effort, selected for cost and availability through the session used to prepare this study. Actual identifiers and settings must be confirmed in execution records. Both are from OpenAI, so their comparison does not establish agreement across providers. Repeated runs of one model are not distinct models or people.

The design allocates three fresh sessions per model, condition and task. Study 2 adds condition C only for real benchmarks. Table~\ref{tab:sessions} gives the total allocation. Three repetitions are a feasibility choice, not a power calculation or evidence of long-run stability.

\begin{table}[t]
\caption{Planned sessions per model. No sessions have been run.}
\label{tab:sessions}
\small
\begin{tabularx}{\columnwidth}{@{}Xccc@{}}
\toprule
Task & A & B & C \\
\midrule
24 fictional cases & 3 & 3 & -- \\
50 real benchmark questions & 3 & 3 & 3 \\
\midrule
Total & 6 & 6 & 3 \\
\bottomrule
\end{tabularx}
\par\smallskip
{\footnotesize Fifteen sessions per model; 30 across both models. A: questions alone. B: explanations and summaries. C: explanations and original passages. B is shared by Studies 1 and 2.}
\end{table}

Execution will interleave conditions under a recorded schedule. Question order and benchmark order will be fixed for matched runs, with changes across repetitions recorded in advance. Two benchmarks in one session may still influence each other. Outputs are therefore summarised by session and benchmark, not treated as independent samples at every answer.

A fresh context alone does not remove workspace access or tools. Before collecting responses, the runner must demonstrate that reviewer sessions do not inherit this research discussion, project records, other answers or reference keys. The available environment and access controls must be recorded. If those conditions cannot be enforced, the protocol must be resolved before execution rather than claiming isolation. The exact inputs, settings, output rules and handling of failures will be frozen after material checking. No such freeze or isolation check is claimed in this draft.

\subsection{Analysis and results}
We will first compare each choice with its reference answer, then examine the reason given for it. This separates answers supported by the supplied information from answers that happen to select the expected choice. Results will show matches, confusion between choices and unsupported conclusions for each model, condition, session and task. Invalid and incomplete outputs remain in the counts. Reading and classifying the reasons requires recorded rules and analyst decisions; matching the choices alone cannot perform that assessment.

\textbf{Results pending.} [Insert attempted, completed and usable session counts; counts for assessable and contested items; fictional-case results; real benchmark results; and examples of interpretation failures. Report the author review status of the reference answers. Do not substitute results from the earlier model study.]

\section{Study 2: Passages and Summaries}
\subsection{Comparison}
The second study turns from the explanations to the benchmark reading material. It reuses condition B and adds condition C, which keeps the questions, explanations, definitions, example and answer format but supplies selected original passages in place of the survey summaries. Reusing B reduces the number of sessions. It also means that the two studies share responses and cannot be treated as independent replications.

The comparison asks what changes when readers encounter a summary rather than selected source passages. Both versions need enough context to interpret the questions, although neither is assumed to contain every relevant fact. We will record their differences in information and use those differences to help explain changes in the answers.

\subsection{Analysis and results}
For each question and benchmark, the analysis will compare the distribution of choices under B and C and the reasons accompanying them. Matched repetition identifiers represent a shared execution schedule, not the same reviewer experiencing both conditions. We will distinguish answers supported by each version, answers exceeding its information, and cases for which the reference remains contested. Examples will link the original passage, summary, changed answer and the specific omission or interpretation involved.

An answer changing from Yes to Not sure can be appropriate if the summary omits a relevant fact. An answer changing from Not sure to No can indicate unsupported certainty, but only after checking the wording and context. Agreement between conditions cannot prove that both are unbiased; both might omit the same qualification.

\textbf{Results pending.} [Insert the answer-distribution comparison, the number of items whose reference differs between packets, unsupported conclusions under each presentation, and traceable examples. Separate model behaviour from the preparers' source audit and retain contested cases.]

\section{Discussion and Limitations}
\subsection{What can be learned from these comparisons}
Together, the studies examine two steps between a review criterion and an answer: understanding what the question asks, and finding what the benchmark documentation says about it. If explanations reduce particular mistakes, that would show how the guidance helps under the tested model conditions. If summaries change answers, checking them against the sources can help identify an omission, an interpretation or a fact made easier to find. These findings could inform the presentation of MESA, while leaving its approachability for people to be examined separately.

RQ1 also sets a boundary on interpretation. Some functions are combined and others omitted, so matching reference answers cannot demonstrate that MESA preserves every source framework. Nor can this design establish superiority over EFPA, OECD or BetterBench as complete review procedures: its comparisons concern MESA presentations. A reviewer may still need technical help, more documents or further analysis before relying on a benchmark for a particular use.

\subsection{Reference answers and preparation bias}
The same research project develops the questionnaire, selects cases and prepares references. This can favour the distinctions already emphasised by MESA. Source links, versions, explicit alternatives and recorded checking make the decisions inspectable but do not remove this dependence. Fictional cases may contain obvious clues and produce ceiling effects; their purpose is to expose selected mistakes, not estimate performance on every benchmark.

Real benchmark references are also limited by source selection. Knowing that a procedure exists somewhere does not mean a participant could infer it from the packet. Some summaries already explain why a particular conclusion is unsupported. Such explanations may teach the intended distinction, so a matching answer can reflect a cue rather than independent reasoning. Reporting the packet as part of the intervention is necessary.

\subsection{Models, scope and prior work}
Only two benchmark descriptions and two models are planned. Model pretraining may contain benchmark facts despite instructions to use supplied text. Repetitions within a model, answers within a session and related fictional cases are dependent. Model families, dates, settings, invalid outputs and tool access must be disclosed. We will not interpret model run duration or completion rate as human effort or understanding.

The earlier CHASE study is preserved separately in \texttt{CHASE-study.md}. It contains 30 completed conversations, 28 usable responses and 32 attempts with Claude Sonnet 5.5. Its comparison used complete EFPA-TRM and general instructions, and only MESA received a teaching example. Those observations are not pooled with the new conditions. Cases used during earlier development are not represented as new tests, and new reference answers cannot make a retrospective analysis prospective.

The relevant ACM SIGSOFT empirical standards guide reporting of the method, comparison conditions, records and limits \cite{sigsoft_empirical_standards}. The IEEE document class concerns presentation, not compliance with research standards. No claim of full compliance or independent validation is made.

\section{Study 3: Learning from Participants}
The existing participant survey provides the next step towards understanding how people experience MESA. This extension will use the survey already prepared, without requiring a second one. Participants encounter familiarity questions, the HLE example, ARC-AGI-2 and an optional DesignQA review. Their experience therefore concerns the questionnaire together with its passages, definitions and example. The later report must identify the packet version and collection procedure used, which may differ from the model study snapshot.

The five feedback items address question clarity, distinguishing answer choices, manageable workload, understanding what to look for and confidence using MESA. Each is reported separately, retaining Cannot judge and missing responses as different categories. These are perceptions, not a demonstrated understanding test or a composite usability score. Optional comments can identify difficulties, but silence does not establish ease of use.

Reference comparisons may describe participant answer patterns on assessable questions using their actual packet. A matching choice does not demonstrate understanding, especially when reasons are optional. Differences from AI answers are descriptive and do not provide a human comparison condition. Familiarity is self-report, submission timestamps do not measure task duration, and the opening and closing confidence items concern different activities. Recruitment, exclusions, analysis timing and repeated responses within participants must be documented in the extension. No human findings are reported here.

\section{Supporting Materials}
The repository is \url{https://github.com/MyLittleFoxxie/MESA-meta-evaluation-standard-for-ai}. The existing \texttt{CHASE-study.md} identifies the prior study. The new preparation record is under \path{surveys/colleague-review/mesa-study/chase-v2/}; it records the packet snapshot, reference preparation status and proposed conditions. The literature review and adaptation maps provide source locations and uncertainties. Material preparation, protocol freeze, execution and analysis are separate stages. Each must be documented before the manuscript describes it as completed.

\section{Conclusion}
MESA starts from a need to understand what an AI benchmark score can support. It brings selected test review concerns into 25 questions, with explanations and five answer choices. The adaptation analysis makes clear what these questions retain and where further investigation is needed. The proposed studies then examine whether explanations help models apply the questions and how summaries affect their answers. Those effects remain to be tested. The existing human survey will later add participants' accounts of using the questionnaire and its supporting material.

\section*{Use of Artificial Intelligence}
OpenAI Codex assisted with the literature synthesis, this manuscript draft and the study preparation records. The new reviewer experiments have not been run. Reference labels, source interpretations and candidate gaps require author checking; no independent screening or coder agreement is claimed. The final report will identify the models that actually produced responses and the people or assistants who analysed them. Mechanically checked counts will remain separate from interpretations of answers and reasons.

\bibliographystyle{IEEEtran}
\bibliography{references}
\end{document}
